\section{Verification Methodology}
\label{sec:validation}

Because the central premise of the mod is that its behavior is not a scripted
approximation, each element's stamping and state-update routines were checked during
development against closed-form analytic solutions -- the exponential RC and RL step
responses, and the memristor's charge-controlled ODE integrated in closed form for a
constant applied voltage -- agreeing to within approximately 1\%, consistent with the
timestep employed ($h=1/20$~s) and the known truncation error of the trapezoidal-integration
companion models. The diode and op-amp were checked the same way: a voltage-follower
configuration reproduced $v_{out}=v_{in}$ to floating-point precision, a non-inverting
amplifier with a $1\,\mathrm{k}\Omega$/$2\,\mathrm{k}\Omega$ feedback divider matched the
textbook gain $1+R_2/R_1=3$ to within roughly one part per million, and the silicon diode
preset converged to a forward drop of approximately $0.69$~V at approximately $4.3$~mA.
Because the solver package has no dependency on any Minecraft class, this class of check runs
as a plain Java program requiring no game instance, which made it practical to re-verify every
element in isolation before integrating it with the rest of the mod. Interface assumptions
against the specific, recently released Minecraft version this mod targets were similarly
checked directly against the decompiled game source rather than against generic recollection
of earlier, better-documented versions. At the time of writing, this verification was
performed ad hoc during development rather than captured as a persisted, automatically
re-run test suite; formalizing it into a checked-in regression suite is planned prior to a
longer-term classroom evaluation (Section~\ref{sec:limitations}).
